% Documentary recovery for Article V; RESULTADO_RECUPERADO.
% Source owner: /Users/ruben/Documents/New project/output/REVISION_EDITORIAL_HMT_MD_20260905/01_FUENTES/INTEGRAL/tree/New project/output/ENSAMBLADO_CAUSAL_RESTAURADO_HMT_MD_20260905/fuente/manuscrito/sections/md/11d_regimenes_cuadraticos_lector_temporal.tex
% Source lines: 4--38, 97--252.
% Only reference labels and heading levels are adapted.
% The recovery matrix preserves the correspondence and dependencies.
\section{Quadratic regimes and temporal readout}
\label{v:rec:chap:regimenes-cuadraticos-tiempo}
\index{quadratic algebra!three regimes}
\index{time crystal!finite HMT realization}
\index{time!TRIT reader}

\subsection{Quadratic algebras of the TRIT and curvature signatures}
\label{v:rec:sec:regimenes-cuadraticos-recibidos}

The primary construction and proof of the three regimes are given in
the construction of the quadratic algebras of the TRIT,
whose defining relation is presented in \eqref{eq:algebras-trit}. For
\(\tau\in\TRIT\), with balanced image
\(\bar\tau\in\{-1,0,+1\}\), we have
\[
 \mathbb A_\tau=\RR[X]/(X^2+\bar\tau),
\]
and, if \(J\) is the class of \(X\),
\[
\begin{array}{ccl}
 \tau=+1&:&J^2=-I,\quad\text{elliptic regime},\\
 \tau=0&:&J^2=0,\quad\text{parabolic regime},\\
 \tau=-1&:&J^2=+I,\quad\text{hyperbolic regime}.
\end{array}
\]
For all three values of \(\tau\), the quotient
\(\mathbb A_\tau=\mathbb R[X]/(X^2+\bar\tau)\) satisfies precisely the three
quadratic identities above. In particular, the internal square root of
\(-1\) is realized as an operator on the elliptic sheet; HMT need neither reject nor
suppress the complex extension. The parabolic algebra contains a nilpotent,
and the hyperbolic algebra decomposes by means of the two idempotents
\((I\pm J)/2\).

These identities classify algebras and transports. They do not, by themselves,
select a spacetime metric or prove that the physical universe occupies one
of the three sheets.

\subsection{Temporal TRIT reader on ({-1,0,+1}): opening, threshold, and return with memory}
\label{v:rec:sec:lector-temporal-trit}

Once an ordered parametrization of the transitions has been chosen, the active
temporal reader is
\[
\begin{aligned}
 +1&\longmapsto
 \text{return, retained memory, and past},\\
 0&\longmapsto
 \text{evaluation threshold and present},\\
 -1&\longmapsto
 \text{bidirectional opening and future}.
\end{aligned}
\]
The reader is \textsc{proved with an explicit starting structure}: its
conclusions follow once the TRIT signature, the order of the
trajectory, and the transported memory have been specified. It does not turn the set
\(\{-1,0,+1\}\) into an equation of motion. The future designates the space
of compatible continuations; survival may select a branch
or a multisection without erasing the genealogy of the others.

The frequency
\[
 \omega_0=\frac{2\pi}{108\,t_0}
\]
assigns a type to the full-return clock through the decision entropy
\(I_{\mathrm{TRIT}}=\log 3\) and the Boltzmann transducer.
It does not, in isolation, establish a stable subharmonic phase. It does, however, enable construction of
the finite reference Hamiltonian below; selection of a robust physical
Hamiltonian requires additional conditions.

\subsection{Unitary dynamics on 108 states: Hamiltonian, propagator, and period}
\label{v:rec:sec:hamiltoniano-reloj-108}

Fix an orientation \(a\in\{\mathrm{pre},\mathrm{ret}\}\) of the
sections constructed in \eqref{v:ant:eq:el-secciones-accion} and write
\(\hbar_{\rm int}=\hbar_a\). The same section is retained in the
Hamiltonian, propagator, and subsequent thermal conversion.

Let
\[
 \mathcal H_{\rm clk}:=\mathbb C^{108},
 \qquad
 \{|j\rangle:j\in\mathbb Z/108\mathbb Z\}
\]
with its canonical Hermitian inner product. Fix
\(\zeta:=\exp(2\pi\mathrm i/108)\), the unitary shift
\[
 U_{\rm clk}|j\rangle=|j+1\rangle
\]
and the Fourier basis
\[
 |\widetilde k\rangle
 :=\frac1{\sqrt{108}}\sum_{j=0}^{107}\zeta^{jk}|j\rangle,
 \qquad 0\leq k\leq107.
\]
With \(t_0\in\mathcal L_T^+\) and
\(\hbar_{\rm int}\in\mathcal L_S^+\), define
\[
 \boxed{
 H_{\rm clk}
 :=\hbar_{\rm int}\omega_0
 \sum_{k=0}^{107}k\,
 |\widetilde k\rangle\langle\widetilde k|,
 \qquad
 \omega_0=\frac{2\pi}{108\,t_0}.}
 \tag{CLK1}\label{v:rec:eq:hamiltoniano-reloj-108}
\]
The operator has energy type
\([H_{\rm clk}]=\mathcal L_S\mathcal L_T^{-1}=\mathcal L_E\).

\begin{theorem}[Exact unitary realization of the \(108\) calendar]
\label{v:rec:thm:reloj-unitario-108}
The operator \(H_{\rm clk}\) is self-adjoint and its one-step propagator
satisfies
\[
 U_{\rm step}
 :=\exp\!\left(-\frac{\mathrm i}{\hbar_{\rm int}}
 H_{\rm clk}t_0\right)
 =U_{\rm clk},
 \qquad
 U_{\rm step}^{108}=I.
 \tag{CLK2}\label{v:rec:eq:propagador-reloj-108}
\]
If
\[
 Z_{\rm clk}|j\rangle=\zeta^j|j\rangle,
\]
then, for any basis state \(|j_0\rangle\),
\[
 \langle Z_{\rm clk}\rangle_{n}
 :=\langle j_0|U_{\rm step}^{-n}Z_{\rm clk}
 U_{\rm step}^n|j_0\rangle
 =\zeta^{j_0+n}.
 \tag{CLK3}\label{v:rec:eq:respuesta-reloj-108}
\]
The sequence has primitive period \(108\).
\end{theorem}

\begin{proof}
Expression \eqref{v:rec:eq:hamiltoniano-reloj-108} is a spectral
decomposition with real eigenvalues, hence \(H_{\rm clk}=H_{\rm clk}^\dagger\).
Moreover,
\[
 U_{\rm clk}|\widetilde k\rangle
 =\zeta^{-k}|\widetilde k\rangle,
\]
whereas the exponential of
\(\hbar_{\rm int}\omega_0 k\) over \(t_0\) gives
\(\exp(-2\pi\mathrm ik/108)=\zeta^{-k}\). The two operators agree on
an orthonormal basis. The remaining identities follow from
\(U_{\rm step}^n|j_0\rangle=|j_0+n\rangle\) and the primitivity of
\(\zeta\).
\end{proof}

The corresponding variational formulation, for a curve
\(\psi\in C^1([t_a,t_b],\mathcal H_{\rm clk})\), unconstrained during the
variation, is
\[
 \boxed{
 L_{\rm clk}(\psi,\dot\psi)
 =\frac{\mathrm i\hbar_{\rm int}}2
 \bigl(
 \langle\psi|\dot\psi\rangle
 -\langle\dot\psi|\psi\rangle
 \bigr)
 -\langle\psi|H_{\rm clk}|\psi\rangle.}
 \tag{CLK4}\label{v:rec:eq:lagrangiano-reloj-108}
\]
The action functional is
\[
 \boxed{
 S_{\rm clk}[\psi]
 :=\int_{t_a}^{t_b}
 L_{\rm clk}(\psi(t),\dot\psi(t))\,\mathrm dt,}
 \tag{CLK5}\label{v:rec:eq:accion-reloj-108}
\]
with \(\delta\psi(t_a)=\delta\psi(t_b)=0\).
\index{Schrödinger!finite-clock equation}
Varying \(\psi\) and \(\bar\psi\) independently yields the
Schrödinger equation for the finite clock
\[
 \boxed{\mathrm i\hbar_{\rm int}\dot\psi=H_{\rm clk}\psi}
 \tag{CLK6}\label{v:rec:eq:schrodinger-reloj-108}
\]
together with its adjoint equation \cite{SchrodingerNobel}. Its propagator sampled
at \(t=nt_0\) is \(\psi_n=U_{\rm step}^{n}\psi_0\). Since \(H_{\rm clk}\) is
self-adjoint, the norm is conserved; normalized initial data remain
normalized. The result is exact with respect to the specified
\(H_{\rm clk},t_0,\hbar_{\rm int}\); it neither selects a physical Hamiltonian
nor introduces an experimental frequency.

The \cref{v:rec:thm:reloj-unitario-108} establishes the unitary dynamics of the calendar
with respect to the specified complex completion, \(t_0\), and \(\hbar_{\rm int}\).
Its realization as a finite periodic system is developed
in Section~\ref{sec:hmtf2-cristal-temporal}, which
specifies the drive, the order observable, the primitive subharmonic response of period
\(108T_{\rm d}\), and finite spectral stability. A
macroscopic material realization constitutes a distinct external identification, not a
gap in the finite result.
